\subsection{Affine proximal oracles}
\label{sec:affine-prox}

Two exact reductions implement the hard-constrained proximal map while
preserving affine feasibility.  The first works directly in the tangent
space; the second enforces the constraint by solving a dual root.

\subsubsection{Tangent-space reduction}

Let
\begin{equation}
  P_{\Tspace}
  :=I-A^\top(AA^\top)^\dagger A
  \label{eq:tangent-projector}
\end{equation}
be the orthogonal projector onto \(\ker A\), where \({}^\dagger\) denotes the
Moore--Penrose inverse.  The orthogonal normal projector is
\(P_{\Tspace^\perp}=I-P_{\Tspace}\).

\begin{lemma}[Normal invariance]
  \label{lem:normal-invariance}
  For every \(z,s\in\R^n\),
  \begin{equation}
    P_\lambda^{\Cset}(z+s)
    =
    P_\lambda^{\Cset}(z+P_{\Tspace}s).
    \label{eq:normal-invariance}
  \end{equation}
\end{lemma}

\begin{proof}
  Fix any \(c\in\Cset\).  Every \(u\in\Cset\) can be written as
  \(u=c+t\) with \(t\in\Tspace\).  Orthogonal decomposition gives
  \[
    \norm{u-(z+s)}^2
    =
    \norm{t-P_{\Tspace}(z+s-c)}^2
    +
    \norm{P_{\Tspace^\perp}(z+s-c)}^2.
  \]
  The second term is independent of \(u\), while
  \[
    P_{\Tspace}(z+s-c)
    =
    P_{\Tspace}(z+P_{\Tspace}s-c).
  \]
  Hence the two proximal objectives differ on \(\Cset\) only by an additive
  constant.  Their unique minimizers are equal.
\end{proof}

Normal invariance also shows that restricting the outer search to the tangent
space loses no polished values:
\begin{equation}
  \inf_{\substack{s\in\R^n\\\norm s\le\delta}}
  f\!\left(P_\lambda^{\Cset}(x+s)\right)
  =
  \inf_{\substack{d\in\Tspace\\\norm d\le\delta}}
  f\!\left(P_\lambda^{\Cset}(x+d)\right).
  \label{eq:fullspace-outer}
\end{equation}
Indeed, \(P_{\Tspace}s\) has norm no larger than \(s\) and gives the same
proximal point, while every tangent \(d\) is already an admissible
  full-space perturbation.  Equation~\eqref{eq:fullspace-outer} shows that
  normal components are computationally redundant as escape directions.

If \(N\in\R^{n\times k}\) has orthonormal columns spanning \(\Tspace\), every
feasible point has the form \(x+Nq\) for a fixed feasible \(x\).  For a tangent
center \(x+N\xi\),
\begin{equation}
  P_\lambda^{\Cset}(x+N\xi)
  =
  x+N\argmin_q
  \left\{
    f(x+Nq)+\frac{1}{2\lambda}\norm{q-\xi}^2
  \right\}.
  \label{eq:nullspace-prox}
\end{equation}
Indeed, \(\norm{N(q-\xi)}=\norm{q-\xi}\), and injectivity of \(N\) makes
  the reduced minimizer unique.  The formula removes the full-space affine
  condition algebraically.
The nullspace formulation is attractive when \(k=n-\operatorname{rank}A\)
is small.  It enforces feasibility exactly in the mathematical formulation
and, in floating-point arithmetic, up to the tolerance of the nullspace
linear algebra, without introducing a penalty parameter.

\subsubsection{Dual-root reduction}

When the unconstrained prox of \(f\) is easier than the nullspace problem, the
affine constraint can instead be enforced through the dual-root construction
of \citet[Section~2.1]{HouTangTohWang2026}, specialized here to the Euclidean root map
\begin{equation}
  H_z(w)
  :=
  A\prox_{\lambda f}(z+A^\top w)-b.
  \label{eq:dual-root}
\end{equation}

\begin{proposition}[Exact affine-prox recovery from a dual root]
  \label{prop:dual-recovery}
  Suppose \(f\) is proper, lower semicontinuous, and
  \(\rho\)-weakly convex, and \(\lambda>0\) with \(\lambda\rho<1\).  If \(H_z(w)=0\), then
  \begin{equation}
    p=\prox_{\lambda f}(z+A^\top w)
    \label{eq:dual-recovery}
  \end{equation}
  is the unique point \(P_\lambda^{\Cset}(z)\).
\end{proposition}

\begin{proof}
  The root condition makes \(p\) feasible.  For every \(u\in\Cset\),
  \begin{align*}
    \norm{u-(z+A^\top w)}^{2}-\norm{u-z}^{2}
    &=
    -2\ip{A(u-z)}{w}+\norm{A^\top w}^{2}\\
    &=
    -2\ip{b-Az}{w}+\norm{A^\top w}^{2},
  \end{align*}
  which is independent of \(u\).  The point \(p\) globally minimizes the
  shifted unconstrained proximal objective and belongs to \(\Cset\);
  therefore it also minimizes that objective over \(\Cset\).  Removing the
  displayed constant shows that it minimizes the constrained proximal
  objective centered at \(z\).  Uniqueness follows from
  \(\lambda\rho<1\).
\end{proof}

\begin{proposition}[Dual-root existence under affine qualification]
  \label{prop:dual-root-existence}
  In addition to the assumptions of \Cref{prop:dual-recovery}, suppose that
  \begin{equation}
    \Cset\cap\operatorname{ri}(\operatorname{dom}f)\ne\varnothing.
    \label{eq:dual-root-cq}
  \end{equation}
  Then, for every center \(z\in\R^n\), the equation \(H_z(w)=0\) has at
  least one solution.
\end{proposition}

\begin{proof}
  Let \(p=P_\lambda^{\Cset}(z)\).  The function
  \[
    u\longmapsto
    \lambda f(u)+\frac12\norm{u-z}^{2}
  \]
  is proper, lower semicontinuous, and
  \((1-\lambda\rho)\)-strongly convex.  The relative-interior condition
  \eqref{eq:dual-root-cq} is the standard qualification for its restriction
  to \(Au=b\).  The exact affine sum rule and convex Fermat rule therefore
  produce a multiplier \(\eta\) such that
  \begin{equation}
    0
    \in
    \lambda\partial f(p)+p-z+A^\top\eta,
    \qquad
    Ap=b
    \label{eq:prox-affine-kkt}
  \end{equation}
  by \cref{eq:affine-sum-rule} and the Fermat rule
  \citep[Theorem~10.1]{RockafellarWets1998}.  Set
  \(w=-\eta\).  The first inclusion in \cref{eq:prox-affine-kkt} is
  equivalent to
  \[
    p=\prox_{\lambda f}(z+A^\top w),
  \]
  because the shifted unconstrained proximal objective is strongly convex.
  The feasibility equation then gives \(H_z(w)=0\).
\end{proof}

The root map has an underlying convex potential.  Define
\[
  h(u):=\lambda f(u)+\frac12\norm u^{2},
  \qquad
  \Psi_z(w):=h^*(z+A^\top w)-\ip{b}{w}.
\]
The function \(h\) is \((1-\lambda\rho)\)-strongly convex; hence \(h^*\) is
differentiable and
\(\nabla h^*(q)=\prox_{\lambda f}(q)\).  It follows that
\begin{equation}
  \nabla\Psi_z(w)=H_z(w),
  \qquad
  \operatorname{Lip}(H_z)
  \le\frac{\norm A^{2}}{1-\lambda\rho}.
  \label{eq:dual-potential}
\end{equation}
Hence \(H_z\) is a monotone Lipschitz gradient field; convexity pertains to the
scalar potential \(\Psi_z\).

Any root in \Cref{prop:dual-recovery} recovers the affine proximal point; the
qualification in \Cref{prop:dual-root-existence} guarantees existence, while
the multiplier \(w\) may be nonunique.  The safe weak-convexity condition
makes the shifted proximal point the unique global solution used by the
escape certificate.

\subsection{Numerical search and rejection}

\subsubsection{Projected envelope search}

Using these proximal oracles, fix the KKT candidate \(x\) and
a trial radius \(\delta>0\), and minimize
\(F(d):=\Moreau_\lambda^{\Cset}(x+d)\) over
\(D_\delta=\Tspace\cap B_\delta(0)\).  Here the optimization variable is the
center perturbation \(d\), not the candidate \(x\).  Finding any
\(d\in D_\delta\) with \(F(d)<f(x)\) is sufficient: the recovered point
\(p=P_\lambda^{\Cset}(x+d)\) then satisfies
\(f(p)\le F(d)<f(x)\) by \cref{eq:recovery-identity}.

Projected gradient descent provides one way of generating trial
perturbations.  At iteration \(k\), a constrained proximal evaluation gives
both \(F(d_k)\) and its gradient.  Since the center \(z_k=x+d_k\) and the
proximal output are feasible, their difference belongs to \(\Tspace\).
Projecting the gradient step back onto \(D_\delta\) retains the radius
constraint:
\begin{align}
  z_k&=x+d_k, &
  p_k&=P_\lambda^{\Cset}(z_k), &
  g_k&=\lambda^{-1}(z_k-p_k),\nonumber\\
  d_{k+1}
  &=
  \Pi_{\Tspace\cap B_\delta(0)}
       (d_k-\alpha_k g_k).
  \label{eq:outer-pg}
\end{align}
To choose a stepsize, note that
\(\Moreau_\lambda^{\Cset}\) is the Moreau envelope of the proper,
lower-semicontinuous, \(\rho\)-weakly convex function
\(f+\ind_{\Cset}\).  Standard weakly convex Moreau calculus therefore gives
\begin{equation}
  L_E
  :=
  \max\left\{
    \frac{1}{\lambda},
    \frac{\rho}{1-\lambda\rho}
  \right\}
  \label{eq:envelope-lipschitz-bound}
\end{equation}
\noindent as a Lipschitz constant for its gradient
\citep[Lemma~2.5]{DavisDrusvyatskiy2022}.
Specialize \(\alpha_k\equiv\alpha\) with \(0<\alpha\le L_E^{-1}\).  The projection
variational inequality and the descent lemma give
\begin{align}
  F(d_{k+1})
  &\le
  F(d_k)-\frac{1}{2\alpha}\norm{d_{k+1}-d_k}^{2}\nonumber\\
  &=
  F(d_k)-\frac{\alpha}{2}
  \norm{\mathcal{G}_{\alpha}(d_k)}^{2},
  \qquad
  \mathcal{G}_{\alpha}(d)
  :=
  \alpha^{-1}
  \left[
    d-\Pi_{D_\delta}(d-\alpha\nabla F(d))
  \right].
  \label{eq:outer-pg-descent}
\end{align}
For example, the first line follows from
\(\ip{\nabla F(d_k)}{d_{k+1}-d_k}
\le-\alpha^{-1}\norm{d_{k+1}-d_k}^{2}\), followed by the
\(L_E\)-smooth descent bound.  The argument is the classical
gradient-projection mechanism \cite{Rosen1960}, here applied to the safe
envelope.  This descent estimate uses an exact proximal point.  The following
bound records the effect of an inexact inner solve.

\begin{lemma}[Inexact projected-envelope step]
  \label{lem:inexact-envelope-step}
  Set \(z=x+d\), \(\mu_\lambda=\lambda^{-1}-\rho>0\), and
  \(Q_z(u)=f(u)+\ind_{\Cset}(u)+\norm{u-z}^2/(2\lambda)\).
  Suppose \(\dist(0,\partial Q_z(\widehat p))\le\varepsilon\), and let
  \(\widehat g=(z-\widehat p)/\lambda\).  Then
  \begin{align}
    \norm{\widehat p-P_\lambda^{\Cset}(z)}
      &\le\frac{\varepsilon}{\mu_\lambda},&
    0\le Q_z(\widehat p)-\Moreau_\lambda^{\Cset}(z)
      &\le\frac{\varepsilon^2}{2\mu_\lambda},\nonumber\\
    \norm{\widehat g-\nabla F(d)}
      &\le\frac{\varepsilon}{1-\lambda\rho}.
      \label{eq:inexact-envelope-gradient}
  \end{align}
  For \(d^+=\Pi_{D_\delta}(d-\alpha\widehat g)\), with
  \(0<\alpha\le L_E^{-1}\),
  \begin{equation}
    F(d^+)\le F(d)-\frac{1}{4\alpha}\norm{d^+-d}^2
      +\frac{\alpha\varepsilon^2}{(1-\lambda\rho)^2}.
    \label{eq:inexact-envelope-descent}
  \end{equation}
\end{lemma}

\begin{proof}
  Apply the strong-convexity argument of \Cref{lem:inexact-safe-prox}
  to \(Q_z\); the gradient estimate follows by dividing the point error
  by \(\lambda\).  Writing \(s=d^+-d\), the projection inequality and
  \(L_E\)-smoothness give
  \[
    F(d^+)\le F(d)-\frac{\norm s^2}{2\alpha}
       +\norm{\widehat g-\nabla F(d)}\norm s.
  \]
  Use \(ab\le a^2/(4\alpha)+\alpha b^2\) to conclude.
\end{proof}

The residual here includes the constraint normal cone.  In particular,
the hypothesis requires a feasible inner point; a small dual-root
feasibility residual alone is not this bound.  A fixed inner tolerance
permits an additive error in descent, rather than the exact monotonicity
of \cref{eq:outer-pg-descent}.

At a KKT candidate, \(P_\lambda^{\Cset}(x)=x\), so initialization at
\(d_0=0\) would leave the iteration at zero.  Nonzero initial perturbations,
possibly with restarts, are therefore needed.  Descent of \(F\) does not by
itself guarantee that its value falls below \(f(x)\); the envelope and
polished values must be tested against that fixed reference value.

A finite collection of the generated perturbations can be used as the set
\(\mathcal D\) in \Cref{alg:safe-finite-escape}.  That algorithm handles the
rejection decision, whereas \cref{eq:outer-pg} supplies trial centers.
The proximal evaluations use the reductions in \Cref{sec:affine-prox};
the rejection test below also addresses inexact inner solves.

\subsubsection{Warm starts and dual-root stopping}

Let \((x,\phi)\) be a KKT pair of \cref{eq:main-problem}.  At the unperturbed
center, the exact dual root is
\begin{equation}
  w_x=-\lambda\phi,
  \label{eq:centered-warm-start}
\end{equation}
because
\[
  x=\prox_{\lambda f}(x-\lambda A^\top\phi).
\]
For a perturbed center \(x+s\), define the normal correction
\begin{equation}
  q_s:=(AA^\top)^\dagger As,
  \qquad
  w^{(0)}(s):=-\lambda\phi-q_s.
  \label{eq:normal-corrected-warm}
\end{equation}
Then
\begin{equation}
  x+s+A^\top w^{(0)}(s)
  =
  x+P_{\Tspace}s-\lambda A^\top\phi.
  \label{eq:warm-input-identity}
\end{equation}
A purely normal change in the center is corrected exactly.  For a tangent
change, \cref{eq:warm-input-identity} gives the identity
\begin{equation}
  H_{x+s}(w^{(0)}(s))
  =
  H_{x+P_{\Tspace}s}(-\lambda\phi).
  \label{eq:warm-residual-identity}
\end{equation}
Since \(H_x(-\lambda\phi)=0\), Lipschitz continuity of the safe
unconstrained prox then gives the initial residual bound
\begin{equation}
  \norm{H_{x+s}(w^{(0)}(s))}
  \le
  \frac{\norm A}{1-\lambda\rho}
  \norm{P_{\Tspace}s}.
  \label{eq:warm-residual-bound}
\end{equation}

\begin{remark}[Scope of the residual bound]
  Equation~\eqref{eq:warm-residual-bound} controls the initial residual.
  Distance to a selected multiplier additionally requires metric regularity
  or strong monotonicity because multipliers may be nonunique and nonsmooth
  proximal maps may generate an interval of dual roots.
\end{remark}

The same potential supplies a safeguard for a basic root iteration.  Let
\[
  L_H:=\frac{\norm A^2}{1-\lambda\rho}
\]
and, when \(L_H>0\), take \(0<\eta\le L_H^{-1}\).  The gradient step
\begin{equation}
  w^+=w-\eta H_z(w)
  \label{eq:dual-gradient-step}
\end{equation}
satisfies
\begin{equation}
  \Psi_z(w^+)
  \le
  \Psi_z(w)-\frac{\eta}{2}\norm{H_z(w)}^2.
  \label{eq:dual-potential-descent}
\end{equation}
Indeed, \cref{eq:dual-potential} and the smooth descent lemma yield the
stronger coefficient
\(\eta(1-\eta L_H/2)\ge\eta/2\).  Backtracking may enforce the displayed
decrease without knowing a sharp Lipschitz constant.

If
\(\widetilde p=\prox_{\lambda f}(z+A^\top w)\), then
\begin{equation}
  \norm{H_z(w)}=\norm{A\widetilde p-b}
  \label{eq:dual-feasibility-residual}
\end{equation}
is exactly the affine feasibility residual of the shifted unconstrained
proximal point.  Together with descent of \(\Psi_z\), it is a natural stopping
diagnostic.  Translating it into objective-value accuracy additionally
requires control of the dual-root error and of any inexact unconstrained prox
computation.

\subsubsection{Finite rejection certificates}

A finite search retains the one-sided validity of the ideal definition.
Given centers \(d_1,\ldots,d_N\in\Tspace\cap B_\delta(0)\), let
\begin{equation}
  \widehat{\Rprox}_{\lambda,\delta}^{(N)}(x)
  :=
  \max_{1\le i\le N}
  \left[
    f(x)-f\!\left(P_\lambda^{\Cset}(x+d_i)\right)
  \right]_+.
  \label{eq:finite-escape}
\end{equation}
By construction,
\begin{equation}
  0
  \le
  \widehat{\Rprox}_{\lambda,\delta}^{(N)}(x)
  \le
  \Rprox_{\lambda,\delta}(x).
  \label{eq:finite-escape-bound}
\end{equation}
Hence a positive sampled value certifies rejection when the inner affine
proximal points are exact or accompanied by rigorous feasible value bounds.
With floating-point values, the comparison should require
\[
  f(u)\le f(x)-\tau_{\rm cert},
\]
where \(\tau_{\rm cert}>0\) dominates a certified bound on the combined
evaluation and inner-solve error.  A zero sampled value does not certify the
absence of an unsampled escape.  Equivalently, if validated objective
enclosures are available, the test is
\[
  f(x)\in[\underline f_x,\overline f_x],\quad
  f(u)\in[\underline f_u,\overline f_u],\quad
  \overline f_u<\underline f_x,
\]
for a certified feasible \(u\).  A small floating-point feasibility
residual alone does not establish the existence of such a feasible witness.

An inexact nullspace computation may still return an explicitly feasible
point \(u=x+Nq\).  A certified inequality \(f(u)<f(x)\) then gives a direct
rejection witness independently of reduced-proximal optimality; identifying
it with the sampled proximal value in \cref{eq:finite-escape} additionally
requires inner optimality.  A dual iterate with \(H_z(w)\ne0\) is infeasible,
so its objective value becomes admissible only after obtaining a certified
feasible point or a problem-specific error enclosure.
\Cref{alg:safe-finite-escape} summarizes the decision logic.

\begin{algorithm}[htbp]
  \caption{Finite proximal-escape search with certified rejection}
  \label{alg:safe-finite-escape}
  \begin{algorithmic}[1]
    \Require KKT candidate \(x\in\Cset\); parameter
      \(\lambda>0\) with \(\lambda\rho<1\); radius \(\delta>0\); finite set
      \(\mathcal D\subset\Tspace\cap B_\delta(0)\)
    \Ensure \textsc{rejected} with a feasible witness, or
      \textsc{unresolved}
    \For{\(d\in\mathcal D\)}
      \State \(z\gets x+d\)
      \State attempt to compute \(P_\lambda^{\Cset}(z)\) by
        \cref{eq:nullspace-prox} or by a certified root of
        \cref{eq:dual-root}
      \If{the inner computation certifies a point \(u\in\Cset\) and
          \(f(u)<f(x)\)}
        \State \Return \textsc{rejected} with witness \(u\)
      \EndIf
    \EndFor
    \State \Return \textsc{unresolved}
  \end{algorithmic}
\end{algorithm}
\FloatBarrier

Directional coverage remains a limitation in high dimension.  If an escape
cone has relative spherical measure \(p\), then \(N\) independent uniform
directions hit it with probability \(1-(1-p)^N\).  Narrow cones can have
rapidly vanishing \(p\) as the tangent dimension grows.  Adaptive directions,
structured subspace searches, or problem-specific deterministic directions
are therefore needed as the tangent dimension grows.

Replacing \(Au=b\) in \cref{eq:cprox} by a penalty
\(\gamma\norm{Au-b}^2/2\) produces a different problem for every finite
\(\gamma\).  An infeasible point may have lower raw objective than every
feasible point and can therefore create a false rejection even when its
constraint residual is small.  Using such an output as an exact feasible
rejection certificate requires an error analysis that converts infeasibility
and objective error into a valid feasible bound.
